\section{Evaluation protocol}

\paragraph{Task lifecycle and correctness.}
Official KernelBench tasks define a reference \code{Model} and expect a
persistent \code{ModelNew}~\citep{kernelbench2025}. We instantiate both from the
same frozen constructor arguments. We restore the random state and move both
modules to the GPU before warmup. Source, initialization, and state hashes record
provenance. Candidates face static policy checks and five seeded input cases.
Verification covers repeated execution, exact output structure, aliasing, input
effects, special values, shape, dtype, and official tolerances. Before timing, a
\code{TorchDispatchMode} audit rejects high level ATen compute and requires an
observed Triton launch~\citep{tillet2019triton}. These checks enforce the task
contract. They are not a security sandbox.

\paragraph{Task selection.}
The study uses official KernelBench L1 commit \code{423217d9}. Before candidate
generation, a memory preflight selected 48 tasks. Selection used a deterministic
round robin across families and lexical order within each family. The cap was
8~GiB of known resident memory. We amended the original target of 50 before the
manifest was frozen. Only 49 of 100 tasks met the cap, and the next memory tier
began at 20~GiB. One feasible task was reserved for a shakedown that did not
enter the results. No candidate or timing field informed this amendment.

\paragraph{Paired screening and fresh confirmation.}
For task $t$, candidate $c$, process $r$, and paired block $b$, define
\begin{equation}
z_{tcrb}=\log(\widetilde T^{\mathrm{eager}}_{trb}/
\widetilde T^{c}_{tcrb}),
\end{equation}
where each $\widetilde T$ is a per launch CUDA event latency on the same input
snapshot~\citep{nvidiaCudaEvents}. One process screens three candidates per task
with 20 randomized paired blocks. We freeze the candidate with the largest
median $z$. Tasks with no valid candidate remain explicit \textsc{invalid} rows.
The selected winner then receives 20 blocks in each of seven fresh processes.
Each process uses new seeds. Four processes run in the first wave and three in
the second, with at least 30 minutes between waves on the same GPU and software
stack. Every interval adapts its launch count to cover at least 1.5~ms. Method
order is randomized. A read and write buffer sized from the reported L2 capacity
perturbs cache state in the same way for both methods. Compile cost is recorded
separately from runtime~\citep{ansel2024pytorch2}. Promotion is measured against
eager execution.

The practical margin is $\delta=0.02$. We fixed this default before timing to
require a meaningful gain over parity. It was not calibrated to the observed T4
noise. Primary outcomes are valid tasks, screening wins, confirmed wins, false
promotions, and median screening to confirmation optimism with a task bootstrap
interval. A secondary process analysis uses a one sided lower bound, a centered
bootstrap test, and Benjamini--Hochberg correction. The executable protocol fixes
the seeds, resampling rules, telemetry, precision, and failure behavior. The
appendix summarizes these choices.

\paragraph{Controls and candidate multiplicity.}
Seven process calibrations compare eager execution with two controls. The first
is an identical wrapper. The second adds a known amount of GPU work. A separate
lifecycle study uses 24 processes to measure reference reconstruction and
transfer with both synchronized host latency and an enclosing CUDA event.
Performance claims are disabled when a required control is absent or fails.

RQ2 uses two deterministic fused8 regimes. The easy regime freezes and confirms
20 candidates per task for $K\leq20$. The parity stress test also freezes 20
variants per task. Each variant preserves the base output and adds calibrated
Triton copy work into unused scratch. Fractional work units copy a prefix of the
first input. Three calibration processes select eight variants per task inside
the declared $[0.98,1.04]$ window. Calibration data do not enter the estimates.
All 32 selected candidates receive screening and seven confirmation processes.
Confirmation starts at least 30 minutes after screening. We use 4,000 resamples
to estimate budgets $K\in\{1,2,3,5,8\}$. The two regimes use different artifacts
and GPUs. Neither is inferred from the KernelBench candidate pool.
